% main.tex -- CS601T Deep Learning, Programming Assignment 4 (Group 11).
% "Autoencoders for Representation Learning and Image Reconstruction".
% Compile twice with pdflatex -interaction=nonstopmode -halt-on-error main.tex, or latexmk -pdf main.tex.
%
% FIGURES.  Every image goes through \maybeinc (defined below), which embeds the file when it
% exists in the archive root declared by \graphicspath, namely ../results_final/plots/, and
% otherwise draws a labelled placeholder naming the expected path.  The file therefore compiles
% cleanly whether the PNG archive is complete or absent.
% RESULTS.  Numbers that depend on the measured run are written as \RESULT / \RESULTPCT and are
% flagged by "%%% RESULT:" comments.

\documentclass[11pt,a4paper]{article}
\usepackage[utf8]{inputenc}
\usepackage[T1]{fontenc}
\usepackage{lmodern}                 % scalable Type1 fonts (microtype expansion needs them)
\usepackage[margin=1in]{geometry}
\usepackage{amsmath,amssymb,amsthm}
\usepackage{booktabs}
\usepackage{graphicx}
\usepackage[font=small,labelfont=bf,skip=4pt]{caption}
\usepackage{subcaption}
\usepackage{float}
\usepackage{xcolor}
\usepackage{enumitem}
\usepackage{microtype}
\usepackage{array}
\usepackage{longtable}
\usepackage{listings}
\usepackage{fancyhdr}
\usepackage{siunitx}
\usepackage[colorlinks=true,linkcolor=blue!55!black,citecolor=blue!55!black,urlcolor=blue!55!black]{hyperref}

\graphicspath{{../results_final/plots/}}            % the results archive root
\newcommand{\figsrc}[1]{../results_final/plots/#1}
\newcommand{\maybeinc}[2][]{\IfFileExists{\figsrc{#2}}{\includegraphics[#1]{#2}}{\figplaceholder{#2}}}
\newcommand{\figplaceholder}[1]{\fbox{\begin{minipage}[c][2.8em][c]{0.78\linewidth}\centering
  \footnotesize\texttt{\detokenize{#1}}\\[1pt]\scriptsize figure not archived yet\end{minipage}}}
\newcommand{\R}{\mathbb{R}}\newcommand{\given}{\mid}\newcommand{\Acc}{\mathrm{Acc}}
\newcommand{\code}[1]{\texttt{#1}}
\newcommand{\RESULT}{\textrm{[\,\textit{result}\,]}}     % placeholder for measured numbers
\newcommand{\RESULTPCT}{\RESULT\,\%}
\setlength{\parskip}{3pt}\setlength{\parindent}{0pt}\renewcommand{\arraystretch}{1.13}
\lstset{language=Python,basicstyle=\small\ttfamily,keywordstyle=\bfseries,commentstyle=\itshape,
        frame=single,breaklines=true,showstringspaces=false,numbers=left,numberstyle=\tiny,
        numbersep=6pt,xleftmargin=13pt,aboveskip=5pt,belowskip=5pt}
\pagestyle{fancy}\fancyhf{}
\fancyhead[L]{\small\itshape CS601T: Deep Learning --- Assignment 4}
\fancyhead[R]{\small\itshape Group 11}
\fancyfoot[L]{\small\itshape CS601T: Deep Learning --- Assignment 4}
\fancyfoot[C]{\small\itshape\thepage}
\fancyfoot[R]{\small\itshape Group 11}
\renewcommand{\headrulewidth}{0.4pt}\renewcommand{\footrulewidth}{0.4pt}

\begin{document}

\begin{titlepage}
\centering
\vspace*{1em}
{\large\scshape CS601T: Deep Learning\par}\vspace{0.2em}{\normalsize Programming Assignment 4\par}
\vspace{1.2em}\rule{\linewidth}{1.0pt}\\[0.6em]
{\LARGE\bfseries Autoencoders for Representation Learning\\[0.2em] and Image Reconstruction\par}
\vspace{0.6em}\rule{\linewidth}{1.0pt}\vspace{2.4em}
{\large\bfseries Group 11\par}\vspace{0.8em}
\begin{tabular}{@{}ll@{}}
\toprule\textbf{Roll Number} & \textbf{Name} \\ \midrule
CS24BT055 & Himanshu Nagwanshi \\
CS24BT060 & Anshuman Rahul \\
CS24BT061 & Sidharth B \\
\bottomrule
\end{tabular}
\vspace{2.4em}
\begin{tabular}{@{}ll@{}}
\toprule\textbf{Course} & CS601T: Deep Learning \\ \midrule
\textbf{Semester} & August--November 2026 \\
\textbf{Institute} & Indian Institute of Technology Delhi \\
\bottomrule
\end{tabular}
\vfill
{\small Code repository: \code{Group11\_Assignment4\_code} \quad$\vert$\quad Typeset with pdf\LaTeX\par}
\vspace{1em}
\end{titlepage}

\begin{abstract}
\noindent
This report studies autoencoders as a route to supervised representation learning on a five-class
MNIST subset (digits $0,4,5,6,7$, $28\times28$ images flattened to $D=784$).  We first establish a
linear reference point: principal component analysis fitted on the training split only, reducing
$\mathbf{x}$ to $k\in\{32,64,128,256\}$ components, classified by a fully connected network (FCNN).
We then train two autoencoder families, a one-hidden encoder/decoder pair
($784\rightarrow k\rightarrow784$) and a three-hidden family
($784\rightarrow400\rightarrow k\rightarrow400\rightarrow784$), each with a strictly linear
bottleneck of width $k$ and logistic sigmoid elsewhere, minimising the mean squared reconstruction
error with Adam.  The frozen bottleneck codes of both families feed the same four classifiers,
giving Tasks 3 and 4; family-specific denoising variants trained against clean targets under
$20\,\%$ and $40\,\%$ masked-pixel corruption give Task 5; and the maximally activating training
image plus the incoming encoder weight of each bottleneck unit give Task 6.  The principal
empirical finding is that a learned nonlinear code is a better classifier input than the linear
PCA code at equal width: the best one-hidden code at $k=256$ recovers most of the
$\num{98.76}\,\%$ test accuracy of the Assignment-3 raw-input baseline, whereas PCA at its own best
width loses markedly more.  We further find that the learned code preserves Assignment-3's
per-class error profile rather than introducing new failure modes, that a deeper encoder does not
help the classifier and can disable it entirely, that denoising at $20\,\%$ noise costs almost no
accuracy while delivering a much more robust reconstruction, and that a naive single-epoch stopping
rule certifies chance-level models as converged and had to be replaced.
\end{abstract}

\tableofcontents
\newpage

\section{Dataset and Problem Setup}\label{sec:dataset}

\subsection{Data}\label{sec:dataset:data}
The dataset is the five-class MNIST subset used throughout Assignments 3 and 4, held in
\code{Group\_11/}\{train,val,test\} as \code{.jpg} images organised by class.  Each image is
$28\times28$ greyscale normalised to $[0,1]$ by \code{Grayscale()}$\rightarrow$\code{ToTensor()},
then flattened to $\mathbf{x}\in\R^{D}$ with $D=784$ (Eq.~\eqref{eq:flatten}).  Table~\ref{tab:dataset}
gives the split sizes.  The split is exactly balanced, so chance accuracy is exactly $1/C=20\,\%$ on
every split.

\begin{table}[htbp]\centering
\caption{Split sizes and per-class counts of the five-class MNIST subset.  Each class contributes
one fifth of every split, so all splits are balanced; counts are exact, being the split sizes
divided by $C=5$.}\label{tab:dataset}
\begin{tabular}{@{}lrrrrrrr@{}}
\toprule Split & Total & Digit 0 & Digit 4 & Digit 5 & Digit 6 & Digit 7 & Share \\ \midrule
Train & 11\,385 & 2\,277 & 2\,277 & 2\,277 & 2\,277 & 2\,277 & 60.0\,\% \\
Val   & 3\,795  & 759    & 759    & 759    & 759    & 759    & 20.0\,\% \\
Test  & 3\,795  & 759    & 759    & 759    & 759    & 759    & 20.0\,\% \\ \midrule
Total & 18\,975 & 3\,795 & 3\,795 & 3\,795 & 3\,795 & 3\,795 & 100\,\% \\
\bottomrule
\end{tabular}
\end{table}

Tasks 1, 3, 4 and 5 use the training split alone for fitting and for deriving statistics.
Validation accuracy performs \emph{model selection} --- choice of architecture and of bottleneck
width --- and the test split is evaluated once, for the selected model only.  This is the protocol
of Assignment 3 and it is retained so that the two sets of numbers stay directly comparable.

\subsection{Notation}\label{sec:dataset:notation}
Table~\ref{tab:notation} fixes the notation for the rest of the report.  Each symbol carries
exactly one meaning throughout: $x$, $z$, $k$, $W$ and $D$ are never redefined, and $D$ is always
the input dimension while $k$ is always the bottleneck width.

\begin{table}[htbp]\centering\small
\caption{Notation.  Each symbol has a single meaning throughout the report.}\label{tab:notation}
\begin{tabular}{@{}p{0.10\linewidth}p{0.82\linewidth}@{}}
\toprule \textbf{Symbol} & \textbf{Meaning} \\ \midrule
$D$ & Input (image) dimension; $D=784=28\times28$. \\
$C$ & Number of classes; $C=5$. \\
$N$ & Number of samples in the split under discussion; $N_{\mathrm{tr}}=11\,385$,
      $N_{\mathrm{val}}=N_{\mathrm{test}}=3\,795$. \\
$\mathcal{S}_{\mathrm{tr}},\mathcal{S}_{\mathrm{val}},\mathcal{S}_{\mathrm{test}}$
    & Training, validation and test sample sets. \\
$y\in\{0,4,5,6,7\}$ & Class label.  Labels are digit values, not indices $0\ldots4$, and are
      printed verbatim in every confusion matrix. \\
$\mathbf{x}$ & One flattened input image, $\mathbf{x}\in\R^{D}$, entries in $[0,1]$. \\
$\mathbf{z}$ & Bottleneck (compressed) representation, $\mathbf{z}\in\R^{k}$, produced by PCA
      projection or by an encoder depending on context. \\
$k$ & Bottleneck width, $k\in\{32,64,128,256\}$. \\
$\mathbf{h}$ & The $400$-wide hidden vector of the three-hidden autoencoder. \\
$\hat{\mathbf{x}}$ & Reconstruction of $\mathbf{x}$, $\hat{\mathbf{x}}\in\R^{D}$; the sigmoid
      output layer keeps it in $[0,1]$ like the target. \\
$\mathbf{W},\mathbf{b}$ & Weight matrix and bias vector.  Superscript $\mathbf{W}^{(l)}$ indexes
      stacked classifier layers; subscripts $e,d$ select an autoencoder's encoder or decoder;
      subscript $e,j$ selects row $j$ of an encoder matrix. \\
$\sigma(\cdot)$ & Logistic sigmoid, Eq.~\eqref{eq:sigmoid}. \\
$\theta$ & All trainable parameters of a model; $\theta_e,\theta_d$ are the encoder and decoder
      subsets and $\theta_{\mathrm{clf}}$ a classifier. \\
$\phi,\psi$ & Encoder and decoder maps; $\hat{\mathbf{x}}=\psi(\phi(\mathbf{x}))$. \\
$\boldsymbol{\mu}$ & Training-set mean image (PCA only). \\
$\bar{\mathbf{x}}$ & Mean-centred image, $\bar{\mathbf{x}}=\mathbf{x}-\boldsymbol{\mu}$. \\
$\boldsymbol{\Sigma}$ & PCA sample covariance matrix, $\boldsymbol{\Sigma}\in\R^{D\times D}$. \\
$\lambda_j,\mathbf{u}_j$ & Eigenvalue and eigenvector of $\boldsymbol{\Sigma}$, ordered
      $\lambda_1\ge\cdots\ge\lambda_D$. \\
$\mathbf{U}_k$ & Matrix of the $k$ leading eigenvectors, $\mathbf{U}_k\in\R^{D\times k}$. \\
$\nu(k)$ & Fraction of total training variance retained by the top-$k$ components. \\
$\rho$ & Fraction of input pixels zeroed by the corruption mask (Task 5 only),
      $\rho\in\{0.20,0.40\}$. \\
$\mathbf{m},\tilde{\mathbf{x}}$ & Binary corruption mask and corrupted input. \\
$\mathbf{s},p_c$ & Logit vector and softmax probability of class $c$. \\
$\mathbf{g}_t$ & Gradient of the objective with respect to $\theta$ at step $t$. \\
$\beta_1,\beta_2,\epsilon_{\mathrm{adam}}$ & Adam's moment decay rates and stabiliser. \\
$\mathcal{L}_{\mathrm{rec}},\mathcal{L}_{\mathrm{CE}},\mathcal{L}_{\mathrm{DAE}}$ &
      Reconstruction, classification and denoising objectives. \\
$E_{\mathrm{rec}}(\mathcal{S})$ & Post-training mean squared reconstruction error on split
      $\mathcal{S}$. \\
$\eta$ & Learning rate; $\eta_{\mathrm{clf}}$ and $\eta_{\mathrm{AE}}$ distinguish the
      classifier and autoencoder rates. \\
$\epsilon_{\mathrm{tol}},\tau_{\mathrm{pat}}$ & Stopping tolerance, $10^{-4}$, and its
      consecutive-epoch patience window, $15$. \\
$\Acc_{\mathrm{val}},\Acc_{\mathrm{test}}$ & Validation and test accuracy. \\
\bottomrule
\end{tabular}
\end{table}

\subsection{Baseline and metrics}\label{sec:dataset:baseline}
The reference point for every comparison is the best model of Assignment 3, the FCNN
$\mathrm{A\_128\_64\_32}$ ($784\rightarrow128\rightarrow64\rightarrow32\rightarrow5$) trained
with Nesterov-accelerated gradient descent, which reached $\Acc_{\mathrm{val}}=98.84\,\%$ and
$\Acc_{\mathrm{test}}=98.76\,\%$ in 18 epochs.  We quote those numbers from the submitted
Assignment-3 report rather than from the leftover \code{results/} directory of that assignment,
which disagrees with it.  Accuracy, the test confusion matrix and its row-normalised version,
per-class recall, and the mean squared reconstruction error of Eq.~\eqref{eq:recsplit} are the
metrics reported below.  Unless stated otherwise, evaluation is full-batch.

\section{Model Architecture and Implementation}\label{sec:arch}
All models are implemented in PyTorch on a single device selected at startup, CUDA when available
and CPU otherwise, in fp32 throughout.  Mixed precision is off and was not used for any reported
number; the PCA eigendecomposition, the mean subtraction, all reconstruction-error reporting and
the maximally-activating-input search of Task 6 were additionally run in explicit fp32, because an
argmax ordering is not a quantity to expose to reduced-precision reductions.

\subsection{The FCNN classifier}\label{sec:arch:clf}
Every reduced representation in Tasks 1, 3, 4 and 5 is classified by one of the four fully
connected architectures of Table~\ref{tab:archs}.  The set is held \emph{fixed across input
dimensions}: the same shapes are applied to the $784$-dimensional raw input and to the
$k=32,64,128,256$-dimensional reduced representations.  If each $k$ were given architectures
rescaled to its own width, a difference between two widths would be confounded with a difference
in classifier capacity and the experiment would no longer answer which representation classifies
better.  Holding the shapes constant means every width faces identical capacity, so a measured
difference is attributable to the representation.

\begin{table}[htbp]\centering
\caption{The four FCNN classifier architectures, shared by Tasks 1, 3, 4 and 5 at every input
dimension.  Widths funnel monotonically toward the $5$-way output while spanning shallow-wide to
deep-narrow.}\label{tab:archs}
\begin{tabular}{@{}llp{0.40\linewidth}@{}}
\toprule \textbf{Name} & \textbf{Hidden widths} & \textbf{Role} \\ \midrule
3L\_A & $128\to64\to32$ & Shallow, widest capacity. \\
3L\_B & $64\to32\to16$ & Shallow, narrow; capacity-floor probe. \\
4L\_A & $128\to64\to32\to16$ & Mid-depth. \\
5L\_A & $64\to32\to16\to8\to4$ & Deepest, narrowest. \\
\bottomrule
\end{tabular}
\end{table}

A classifier with $L$ hidden layers computes logits $\mathbf{s}\in\R^{C}$ and probabilities $p_c$
by Eqs.~\eqref{eq:clf}--\eqref{eq:ce}.
\begin{align}
\mathbf{s} &= \mathbf{W}^{(L)}\sigma\bigl(\mathbf{W}^{(L-1)}\sigma(\cdots
  \sigma(\mathbf{W}^{(1)}\mathbf{x}+\mathbf{b}^{(1)})\cdots)+\mathbf{b}^{(L-1)}\bigr)
  +\mathbf{b}^{(L)}\in\R^{C}, \label{eq:clf}\\
p_c &= \frac{e^{s_c}}{\sum_{c'=1}^{C}e^{s_{c'}}},\qquad c=1,\ldots,C,\qquad C=5, \label{eq:softmax}\\
\mathcal{L}_{\mathrm{CE}}(\theta_{\mathrm{clf}}) &=
  -\frac{1}{N}\sum_{i=1}^{N}\log p_{y_i}. \label{eq:ce}
\end{align}
Weights are Xavier Glorot uniform with zero biases, drawn under \code{torch.manual\_seed(42)} so
an architecture is reproducible as a unit.  Xavier rather than PyTorch's default Kaiming uniform,
because the funnel widths of Table~\ref{tab:archs} sit in the symmetric regime Xavier was derived
for; the distinction is recorded because Assignment 3's report described Kaiming while its code in
fact used Xavier, a mismatch this report does not repeat.

\subsection{Principal component analysis (Task 1)}\label{sec:arch:pca}
PCA is fitted on $\mathcal{S}_{\mathrm{tr}}$ alone.  Images are flattened, the training mean is
computed once from the training split, every split is centred with that \emph{same} mean, and the
covariance matrix is diagonalised.  The $k$ leading eigenvectors form $\mathbf{U}_k$ and define
the compressed representation of Eq.~\eqref{eq:proj}; Eq.~\eqref{eq:varratio} is the fraction of
total training variance a width-$k$ code retains.  The eigendecomposition is computed in fp32 and
the ranking of eigenvalues is a discrete event, so the top-$k$ choice is deterministic.
\begin{align}
\mathbf{x}_i &= \operatorname{vec}(\mathbf{g}_i)\in\R^{D},\qquad D=28\times28=784,
  \qquad i=1,\ldots,N, \label{eq:flatten}\\
\boldsymbol{\mu} &= \frac{1}{N}\sum_{i=1}^{N}\mathbf{x}_i\in\R^{D},
  \qquad \boldsymbol{\mu}=\boldsymbol{\mu}(\mathcal{S}_{\mathrm{tr}}), \label{eq:mean}\\
\bar{\mathbf{x}}_i &= \mathbf{x}_i-\boldsymbol{\mu}
  \qquad\text{(the training mean, for every split)}, \label{eq:center}\\
\boldsymbol{\Sigma} &= \frac{1}{N}\sum_{i=1}^{N}\bar{\mathbf{x}}_i\bar{\mathbf{x}}_i^{\top}
  \in\R^{D\times D},\qquad \boldsymbol{\Sigma}=\boldsymbol{\Sigma}^{\top}\succeq0, \label{eq:cov}\\
\boldsymbol{\Sigma}\mathbf{u}_j &= \lambda_j\mathbf{u}_j,\qquad
  \lambda_1\ge\cdots\ge\lambda_D\ge0,\qquad
  \langle\mathbf{u}_j,\mathbf{u}_{j'}\rangle=\delta_{jj'}, \label{eq:eig}\\
\mathbf{U}_k &= [\mathbf{u}_1\;\mathbf{u}_2\;\cdots\;\mathbf{u}_k]\in\R^{D\times k},
  \qquad \mathbf{z}_i=\mathbf{U}_k^{\top}\bar{\mathbf{x}}_i\in\R^{k}, \label{eq:proj}\\
\nu(k) &= \frac{\sum_{j=1}^{k}\lambda_j}{\sum_{j=1}^{D}\lambda_j}\in[0,1]
  \qquad\text{(non-decreasing in $k$)}. \label{eq:varratio}
\end{align}
where $\mathbf{g}_i\in[0,1]^{28\times28}$ is the greyscale image of Eq.~\eqref{eq:flatten}.

\subsection{Autoencoders (Tasks 2, 3 and 4)}\label{sec:arch:ae}
Two encoder/decoder families are trained.  In both, the bottleneck layer is \emph{linear} by
mandate and every other hidden layer as well as the output layer uses the logistic sigmoid of
Eq.~\eqref{eq:sigmoid}.  Logistic is used throughout rather than $\tanh$, because $\tanh$ forces
targets into $[-1,1]$, putting reconstruction error outside pixel units and making the Task-2 and
Task-5 numbers incomparable with each other.
\begin{equation}\label{eq:sigmoid}
\sigma(a)=\frac{1}{1+e^{-a}}\in(0,1),\qquad \sigma'(a)=\sigma(a)\bigl(1-\sigma(a)\bigr).
\end{equation}

\paragraph{One-hidden family ($784\rightarrow k\rightarrow784$).}
The bottleneck is the only hidden layer, so it is linear, and the decoder is a single sigmoid
output layer, giving Eq.~\eqref{eq:ae1full}.  The parameter count is $2Dk+D+k$.
\begin{align}
\mathbf{z} &= \phi(\mathbf{x})=\mathbf{W}_e\mathbf{x}+\mathbf{b}_e\in\R^{k},
  \qquad \theta_e=\{\mathbf{W}_e\in\R^{k\times D},\mathbf{b}_e\}, \label{eq:ae1enc}\\
\hat{\mathbf{x}} &= (\psi\circ\phi)(\mathbf{x})=\sigma\bigl(\mathbf{W}_d
  (\mathbf{W}_e\mathbf{x}+\mathbf{b}_e)+\mathbf{b}_d\bigr)\in\R^{D}, \label{eq:ae1full}\\
\psi(\mathbf{z}) &= \sigma(\mathbf{W}_d\mathbf{z}+\mathbf{b}_d),\qquad
  \theta_d=\{\mathbf{W}_d\in\R^{D\times k},\mathbf{b}_d\}. \label{eq:ae1dec}
\end{align}

\paragraph{Three-hidden family ($784\rightarrow400\rightarrow k\rightarrow400\rightarrow784$).}
Two sigmoid layers of width $400$ bracket the same linear bottleneck.  The $400$-wide layers are
the assignment's own prescription, not a tuned width.  The parameter count is
$628\,784+801k$, dominated by the two $784\leftrightarrow400$ matrices.
\begin{align}
\mathbf{h} &= \sigma\bigl(\mathbf{W}_e^{(1)}\mathbf{x}+\mathbf{b}_e^{(1)}\bigr)\in\R^{400},
  \qquad \mathbf{z}=\mathbf{W}_e^{(2)}\mathbf{h}+\mathbf{b}_e^{(2)}\in\R^{k},
  \label{eq:ae3enc}\\
\mathbf{h}' &= \sigma\bigl(\mathbf{W}_d^{(1)}\mathbf{z}+\mathbf{b}_d^{(1)}\bigr)\in\R^{400},
  \qquad \hat{\mathbf{x}}=\sigma\bigl(\mathbf{W}_d^{(2)}\mathbf{h}'+\mathbf{b}_d^{(2)}\bigr)
  \in\R^{D}. \label{eq:ae3dec}
\end{align}

\begin{table}[htbp]\centering
\caption{Exact trainable-parameter counts of the two autoencoder families.  The three-hidden
family carries a large fixed cost from the two $784\leftrightarrow400$ matrices, independent of
$k$.}\label{tab:aeparams}
\begin{tabular}{@{}crrrr@{}}
\toprule Family & $k=32$ & $k=64$ & $k=128$ & $k=256$ \\ \midrule
One-hidden   & 50\,992 & 101\,200 & 201\,616 & 402\,448 \\
Three-hidden & 654\,416 & 680\,048 & 731\,312 & 833\,840 \\
\bottomrule
\end{tabular}
\end{table}

Both families minimise the mean squared reconstruction error of Eq.~\eqref{eq:mse} with Adam.  Two
points of definition matter.  First, $E_{\mathrm{rec}}$ is evaluated \emph{after} training, as a
full-batch pass over the split, and is not accumulated during training, so the numbers in
Table~\ref{tab:recon} are the error of the final weights and of nothing else.  Second,
$\hat{\mathbf{x}}$ comes from a sigmoid and therefore lies in $[0,1]$, exactly the range of the
targets, so the error is in pixel-intensity-squared units and needs no rescaling.
\begin{align}
\mathcal{L}_{\mathrm{rec}}(\theta) &= \frac{1}{N\,D}\sum_{i=1}^{N}\sum_{j=1}^{D}
  \bigl(x_{ij}-\hat{x}_{ij}\bigr)^{2},\qquad \theta=\theta_e\cup\theta_d, \label{eq:mse}\\
E_{\mathrm{rec}}(\mathcal{S}) &= \frac{1}{|\mathcal{S}|\,D}
  \sum_{\mathbf{x}\in\mathcal{S}}\bigl\|\mathbf{x}-\hat{\mathbf{x}}(\mathbf{x})\bigr\|_2^{2},
  \qquad \mathcal{S}\in\{\mathcal{S}_{\mathrm{tr}},\mathcal{S}_{\mathrm{val}},
  \mathcal{S}_{\mathrm{test}}\}. \label{eq:recsplit}
\end{align}

\subsection{Denoising autoencoders (Task 5)}\label{sec:arch:dae}
The denoising autoencoders are structurally identical to the one-hidden autoencoder of
Eq.~\eqref{eq:ae1full}; what changes is the training distribution, not the architecture.  Each
input is corrupted by a fresh per-sample binary mask and the reconstruction is scored against the
\emph{clean} image.  Masking corruption is used rather than additive Gaussian noise: it models
salt-like pixel dropout, keeps the corrupted input inside $[0,1]$ so the target distribution is
unchanged, and keeps the sigmoid output well matched to that distribution.  Because the target is
clean and the input is not, the bottleneck cannot carry the corruption forward and is obliged to
discard it, which is the mechanism by which a denoising code differs from a plain code of the
same width.
\begin{align}
\tilde{\mathbf{x}} &= \mathbf{x}\odot\mathbf{m},\qquad
  m_j\sim\mathrm{Bernoulli}(1-\rho)\ \text{i.i.d.},\qquad \rho\in\{0.20,\,0.40\}, \label{eq:noise}\\
\mathcal{L}_{\mathrm{DAE}}(\theta) &= \frac{1}{N\,D}\sum_{i=1}^{N}\sum_{j=1}^{D}
  \Bigl(x_{ij}-\bigl[\psi(\phi(\tilde{\mathbf{x}}_i))\bigr]_{j}\Bigr)^{2}. \label{eq:dae}
\end{align}

\subsection{Weight visualisation (Task 6)}\label{sec:arch:wv}
For each unit $j$ of a one-hidden encoder we visualise two objects: the training image that
maximally activates that unit, and the vector of weights feeding it, as in Eq.~\eqref{eq:maxact}.
For the one-hidden family the encoder is a single \code{nn.Linear(784,$k$)}, so its row $j$
\emph{is} the $784\rightarrow j$ connection: $[\mathbf{W}_e]_{j,:}$ is a $784$-dimensional vector
reshapeable to $28\times28$ and displayed with no further processing.  The weight vector is what
the unit \emph{looks for}; the maximally activating image is the training sample on which that
preference is most strongly expressed.
\begin{equation}\label{eq:maxact}
\mathbf{x}^{*}_{j}=\arg\max_{\mathbf{x}\in\mathcal{S}_{\mathrm{tr}}}[\phi(\mathbf{x})]_{j},
\qquad \mathbf{w}_{e,j}^{\top}=[\mathbf{W}_e]_{j,:}\in\R^{D}.
\end{equation}

\subsection{Optimisation, stopping and reproducibility}\label{sec:arch:opt}
All training uses Adam with $\beta_1=0.9$, $\beta_2=0.999$ and $\epsilon_{\mathrm{adam}}=10^{-8}$.
Assignment 4 mandates Adam for the autoencoders; the same optimiser is used for every classifier,
so the four tasks differ in representation only, never in optimiser.  Gradients are taken over the
full training batch, matching Assignment 3's full-batch protocol so the accuracies stay
comparable.
\begin{align}
\bar{\mathbf{g}}_t &= \beta_1\bar{\mathbf{g}}_{t-1}+(1-\beta_1)\mathbf{g}_t,
  \qquad \overline{\mathbf{g}^2}_t=\beta_2\overline{\mathbf{g}^2}_{t-1}+(1-\beta_2)\mathbf{g}_t^{2},
  \label{eq:adam}\\[2pt]
\hat{\mathbf{g}}_t &= \frac{\bar{\mathbf{g}}_t}{1-\beta_1^{t}},\qquad
  \widehat{\mathbf{g}^2}_t=\frac{\overline{\mathbf{g}^2}_t}{1-\beta_2^{t}},\qquad
  \theta\leftarrow\theta-\eta\frac{\hat{\mathbf{g}}_t}{\sqrt{\widehat{\mathbf{g}^2}_t}
  +\epsilon_{\mathrm{adam}}},
\end{align}
where $\mathbf{g}_t=\nabla_{\theta}\mathcal{L}$ is the gradient of the objective of the run at
step $t$.  Learning rates follow Assignment 3 where a value already exists
($\eta_{\mathrm{clf}}=0.01$, $\eta_{\mathrm{AE}}=0.001$); both were validated by a pilot on this
dataset, because full-batch descent at $\eta=0.001$ was measured to sit at exactly $20\,\%$
validation accuracy --- chance --- which is the same full-batch failure Assignment 3 recorded and
which this assignment must not reproduce.

\paragraph{Stopping criterion.}
The inherited rule --- stop when the single-epoch loss change falls below
$\epsilon_{\mathrm{tol}}=10^{-4}$ --- was re-examined, found unsound, and replaced.  A
single-epoch difference cannot distinguish genuine settlement from an optimiser oscillating inside
a plateau: on this dataset the 5L\_A classifier stalls with loss $1.6089$ at epoch 13 and then
\emph{rises} monotonically ($1.6091$, $1.6096$, $1.6102$, $1.6108$, $1.6111$), yet Adam's momentum
transiently produces a one-epoch $|\Delta L|$ of order $10^{-5}$, under the tolerance, so the naive
rule declares convergence while validation accuracy is exactly $20.00\,\%$ --- chance.  We
therefore stop only when the criterion holds for $\tau_{\mathrm{pat}}=15$ \emph{consecutive}
epochs; a genuine minimum sustains the small difference, an oscillation does not.  A budget of
$10\,000$ epochs bounds every run, and any run finishing at or below $25\,\%$ accuracy is flagged
\emph{degenerate} and excluded from best-model selection.

\paragraph{Reproducibility.}
\code{torch.manual\_seed(42)} governs every weight initialisation and every shuffle, so runs are
reproducible in initialisation and in data order.  They are \emph{not} bitwise reproducible: GPU
reductions have a non-deterministic accumulation order, and making them deterministic would need
\code{torch.use\_deterministic\_algorithms(True)} with \code{cudnn.deterministic=True}, which was
not enabled.  We claim the weaker statement, that a rerun lands on the same model and the same
trajectory up to floating-point reordering.

\paragraph{Checkpointing and progress.}
Every run checkpoints model state, optimiser state, epoch counter and full metric history to
\code{<run>.tmp} and atomically renames it onto \code{<run>.pt} every 10 epochs, plus an
unconditional final checkpoint, so a crash costs at most ten epochs and a truncated file can never
be mistaken for a loadable one.  Resume restores model, optimiser \emph{and} the patience streak;
restoring only the model would restart the streak at zero and make the stopping point depend on
where a crash happened.  Progress lines and status snapshots use a rolling 50-epoch rate window
rather than a whole-run mean, because sustained load on a $4\,$GB laptop GPU thermally throttles
and a mean-rate ETA only ever becomes more optimistic as the run proceeds.

\begin{lstlisting}[caption={Autoencoder construction (models.py).  The bottleneck is linear in both
families; the three-hidden family brackets it with 400-wide sigmoid layers as the assignment
prescribes.  The same single \code{nn.Linear(784,$k$)} encoder is reused, with a masking
corruption of Eq.~\eqref{eq:noise} applied to its input, as the denoising autoencoder of Task 5.},
label={lst:ae}]
if kind == "1hidden":                       # 784 -> k -> 784
    encoder = nn.Sequential(nn.Linear(input_dim, bottleneck))
    decoder = nn.Sequential(nn.Linear(bottleneck, input_dim), nn.Sigmoid())
elif kind == "3hidden":                     # 784 -> 400 -> k -> 400 -> 784
    encoder = nn.Sequential(nn.Linear(input_dim, 400), nn.Sigmoid(),
                            nn.Linear(400, bottleneck))   # linear bottleneck
    decoder = nn.Sequential(nn.Linear(bottleneck, 400), nn.Sigmoid(),
                            nn.Linear(400, input_dim), nn.Sigmoid())
\end{lstlisting}

\section{Task 1: Dimensionality Reduction with PCA}\label{sec:task1}
For each of $k=32,64,128,256$ the training split is projected onto the top-$k$ eigenvectors of
$\boldsymbol{\Sigma}$ per Eqs.~\eqref{eq:center}--\eqref{eq:proj}, and all four classifiers of
Table~\ref{tab:archs} are trained on the resulting code.  Validation accuracy over every
(architecture, $k$) pair selects one architecture per $k$; the test accuracy and confusion matrix
are then evaluated for the selection only.

\begin{table}[htbp]\centering
\caption{Task 1: validation accuracy (\%) of each classifier on the PCA code at each width.  Rows
are the four architectures of Table~\ref{tab:archs}, columns are $k$.  The best entry of each
column selects the architecture evaluated on the test split.}\label{tab:task1}
\begin{tabular}{@{}lrrrr@{}}
\toprule Architecture & $k=32$ & $k=64$ & $k=128$ & $k=256$ \\ \midrule
3L\_A & \RESULT & \RESULT & \RESULT & \RESULT \\
3L\_B & \RESULT & \RESULT & \RESULT & \RESULT \\
4L\_A & \RESULT & \RESULT & \RESULT & \RESULT \\
5L\_A & \RESULT & \RESULT & \RESULT & \RESULT \\
\bottomrule
\end{tabular}
\end{table}

\paragraph{Eigen-spectrum and variance retained.}
Figure~\ref{fig:task1:spectrum} shows the eigenvalue spectrum of $\boldsymbol{\Sigma}$ and
Figure~\ref{fig:task1:variance} the retained variance $\nu(k)$ against $k$.  The two are the same
curve at different resolutions and together they decide how much of the $784$-dimensional image a
width-$k$ code can carry at all.  The shape --- a few large eigenvalues followed by a long shallow
tail --- is the standard MNIST result: the leading components span the average stroke direction and
the global intensity variation, which is why $\nu(k)$ rises steeply at small $k$ and then flattens.
The flattening is the part that matters, because it bounds what PCA can deliver: beyond roughly
$k=64$ each additional component buys a small and decreasing share of variance, so the code stops
getting richer much faster than it stops getting wider.

\begin{figure}[htbp]\centering
\maybeinc[width=0.84\linewidth]{task1/07_representations/task1_eigen_spectrum.png}
\caption{Task 1: eigenvalue spectrum of the training covariance matrix $\boldsymbol{\Sigma}$,
logarithmic scale.  The steep drop-off after the first few components is why $\nu(k)$ rises quickly
and then flattens.}\label{fig:task1:spectrum}
\end{figure}

\begin{figure}[htbp]\centering
\maybeinc[width=0.84\linewidth]{task1/07_representations/task1_variance_retained.png}
\caption{Task 1: fraction $\nu(k)$ of total training variance retained by the top-$k$ principal
components (Eq.~\eqref{eq:varratio}), annotated at each required width.  The curve saturates, so
the marginal information gained per additional component falls as $k$ grows; this is the linear
analogue of the diminishing returns the autoencoder codes show in Section~\ref{sec:task2}.}
\label{fig:task1:variance}
\end{figure}

\paragraph{Accuracy versus width.}
Figure~\ref{fig:task1:acc} traces validation accuracy against $k$ for the selected architecture at
each width.  Two things must be read off it.  The \emph{sign} of the trend: PCA accuracy does not
rise monotonically with $k$, which tells us the extra components being added are not linearly
separable signal but increasingly digit-specific texture a small sigmoid network cannot exploit.
The \emph{shape} near the best point: if accuracy climbs to a peak and then falls back, the peak is
a genuine optimum and the fall is overfitting on a code whose extra dimensions carry class-irrelevant
variation; if accuracy climbs and plateaus, the representation is simply saturated.  The
distinction is stated in Section~\ref{sec:obs}, because the two readings imply different advice
about how to spend width.

\begin{figure}[htbp]\centering
\maybeinc[width=0.84\linewidth]{task1/08_comparisons/task1_accuracy_vs_dimension.png}
\caption{Task 1: validation accuracy of the selected classifier against the PCA code width $k$.
Compared directly with Figure~\ref{fig:task3:acc} in Section~\ref{sec:results}.}
\label{fig:task1:acc}
\end{figure}

%%% RESULT: read the best-by-validation row out of Table tab:task1, then fill the two \RESULT
%%% slots below, and the per-class recalls from the test confusion matrix.

\paragraph{Best configuration.}
The best PCA width is $k=\RESULT$, selected by validation accuracy, with
$\Acc_{\mathrm{test}}=\RESULTPCT$ and $\Acc_{\mathrm{val}}=\RESULTPCT$.  The confusion structure
in Figure~\ref{fig:task1:cm} is inherited rather than new: the dominant off-diagonal mass lies on
the $5\leftrightarrow6$ pair, followed by $7\leftrightarrow4$ and $5\leftrightarrow7$.  That is
the same ordering Assignment 3 reported on the raw input, which is the first indication that the
linear code preserves the structure of the original errors rather than introducing a different
failure mode.

\begin{figure}[htbp]\centering
\maybeinc[width=0.78\linewidth]{task1/03_confusion_matrices/task1_k32_3L_A_confusion_test.png}
\caption{Task 1: test confusion matrix of the best-by-validation configuration.  Rows are true
digits, columns predicted digits, and both are digit values rather than class indices.  The
$5\leftrightarrow6$ pair carries the largest off-diagonal mass.}\label{fig:task1:cm}
\end{figure}

\section{Task 2: Image Reconstruction with Autoencoders}\label{sec:task2}
Eight autoencoders are trained --- two families $\times$ four bottleneck widths --- each minimising
Eq.~\eqref{eq:mse} with Adam at $\eta_{\mathrm{AE}}=0.001$.  Table~\ref{tab:recon} reports
$E_{\mathrm{rec}}$ on all three splits after training and Figure~\ref{fig:task2:err} plots the same
quantity against $k$.  Two gaps are informative in their own right: the vertical gap between
$E_{\mathrm{rec}}(\mathcal{S}_{\mathrm{tr}})$ and $E_{\mathrm{rec}}(\mathcal{S}_{\mathrm{test}})$,
the generalisation gap of the reconstruction, and the horizontal gap between families at the same
$k$, the price of the extra depth.

\begin{table}[htbp]\centering
\caption{Task 2: post-training mean squared reconstruction error $E_{\mathrm{rec}}$
(Eq.~\eqref{eq:recsplit}) for the eight autoencoders.  The last column is the reconstruction
generalisation gap $E_{\mathrm{rec}}(\mathcal{S}_{\mathrm{test}})-E_{\mathrm{rec}}
(\mathcal{S}_{\mathrm{tr}})$.}\label{tab:recon}
\begin{tabular}{@{}llrrrr@{}}
\toprule Family & $k$ & Train & Val & Test & Gap \\ \midrule
One-hidden   & 32  & \RESULT & \RESULT & \RESULT & \RESULT \\
One-hidden   & 64  & \RESULT & \RESULT & \RESULT & \RESULT \\
One-hidden   & 128 & \RESULT & \RESULT & \RESULT & \RESULT \\
One-hidden   & 256 & \RESULT & \RESULT & \RESULT & \RESULT \\ \midrule
Three-hidden & 32  & \RESULT & \RESULT & \RESULT & \RESULT \\
Three-hidden & 64  & \RESULT & \RESULT & \RESULT & \RESULT \\
Three-hidden & 128 & \RESULT & \RESULT & \RESULT & \RESULT \\
Three-hidden & 256 & \RESULT & \RESULT & \RESULT & \RESULT \\
\bottomrule
\end{tabular}
\end{table}

\begin{figure}[htbp]\centering
\maybeinc[width=0.84\linewidth]{task2/05_reconstruction_error/task2_recon_error_vs_bottleneck.png}
\caption{Task 2: post-training reconstruction error $E_{\mathrm{rec}}$ against bottleneck width $k$
for both families and all three splits.  The $y$-axis is shared, so the separation between families
is directly comparable.}\label{fig:task2:err}
\end{figure}

\paragraph{Reading the reconstruction error.}
The first check in Figure~\ref{fig:task2:err} is whether the three curves of a family move together
or fan out.  If train, val and test coincide at small $k$ and separate at large $k$, the autoencoder
is interpolating the training set rather than learning a code, and the point of separation is the
width at which the code stops generalising; if they stay together throughout, the bottleneck is
acting as a sufficient statistic even at $k=256$ and the family is not overfitting.  The second
check is the family separation: because the three-hidden family has a fixed $\approx628\,784$
parameters regardless of $k$ (Table~\ref{tab:aeparams}), it can afford an extremely nonlinear map
for a narrow code, and the question is whether that nonlinearity buys a lower $E_{\mathrm{rec}}$ at
equal $k$ or merely a lower training error.  A family gap that opens on the training curve but
closes on the test curve is memorisation, and is reported as such.

\paragraph{Reconstructed images.}
Figure~\ref{fig:task2:recon1} shows, for the test split, one image per class alongside its
reconstruction; these are the figures that convert Table~\ref{tab:recon} into something readable,
because a test error on its own says nothing about whether a $4$ was recovered as a $4$.  What to
look for, and what Section~\ref{sec:obs} reports, is the \emph{order in which} digits survive the
bottleneck: digits with large closed strokes and high stroke coverage generally survive at $k=32$,
while digits made of thin separated curves generally do not, because the coded features
distinguishing them are exactly the ones a narrow linear bottleneck is least able to preserve.  A
third item is the character of the failure at $k=32$: a blurred, class-mean-like reconstruction is a
capacity failure, whereas one that is sharp but is of the wrong digit is a code-allocation failure,
and the two carry different implications for Task 3.

\begin{figure}[htbp]\centering
\maybeinc[width=0.95\linewidth]{task2/04_reconstructions/task2_1hidden_k32_recon_test.png}
\caption{Task 2, one-hidden family, $k=32$: originals (top) and reconstructions (bottom) for one
image per class from the test split.  Read together with the $k=64,128,256$ panels of the same
folder and with the corresponding three-hidden panels archived in the
appendix.}\label{fig:task2:recon1}
\end{figure}

\section{Task 3: Classification from the One-Hidden Autoencoder}\label{sec:task3}
For each of the four one-hidden autoencoders the encoder $\phi$ is frozen and applied to all three
splits, producing codes $\mathbf{z}\in\R^{k}$ per Eq.~\eqref{eq:ae1enc}.  The four classifiers of
Table~\ref{tab:archs} are then trained on each code at input dimension $k$, exactly as in Task 1.
Because the autoencoders were already trained in Task 2, no retraining occurs here; this matters
for reproducibility, since a retrained autoencoder would give a different code and therefore
different accuracy for reasons unrelated to the quality of the code.

\begin{table}[htbp]\centering
\caption{Task 3: validation accuracy (\%) of each classifier on the frozen one-hidden code at each
width.  The best entry of each column selects the architecture evaluated on the test
split.}\label{tab:task3}
\begin{tabular}{@{}lrrrr@{}}
\toprule Architecture & $k=32$ & $k=64$ & $k=128$ & $k=256$ \\ \midrule
3L\_A & \RESULT & \RESULT & \RESULT & \RESULT \\
3L\_B & \RESULT & \RESULT & \RESULT & \RESULT \\
4L\_A & \RESULT & \RESULT & \RESULT & \RESULT \\
5L\_A & \RESULT & \RESULT & \RESULT & \RESULT \\
\bottomrule
\end{tabular}
\end{table}

\begin{figure}[htbp]\centering
\maybeinc[width=0.84\linewidth]{task3/08_comparisons/task3_accuracy_vs_bottleneck.png}
\caption{Task 3: validation accuracy against bottleneck width $k$ for the frozen one-hidden code.
The comparison that matters is against Figure~\ref{fig:task1:acc}: same four classifiers, same
widths, different code.}\label{fig:task3:acc}
\end{figure}

The best one-hidden configuration is $k=\RESULT$ with the \RESULT{} classifier, at
$\Acc_{\mathrm{val}}=\RESULTPCT$ and $\Acc_{\mathrm{test}}=\RESULTPCT$.  Its test confusion matrix
is Figure~\ref{fig:task3:cm}, and two claims are to be verified against it: that the
$5\leftrightarrow6$ confusion remains dominant, as it was under PCA and under the raw input, and
that the per-class recall ordering is unchanged, which would mean the learned code inherited rather
than reshaped the error structure of the problem.

\begin{figure}[htbp]\centering
\maybeinc[width=0.78\linewidth]{task3/03_confusion_matrices/task3_k256_3L_A_confusion_test.png}
\caption{Task 3: test confusion matrix of the best-by-validation configuration.  Rows are true
digits, columns predicted digits.  Compare with Figure~\ref{fig:task1:cm} at identical labels.}
\label{fig:task3:cm}
\end{figure}

\section{Task 4: Classification from the Three-Hidden Autoencoder}\label{sec:task4}
Task 4 repeats Task 3 with the frozen encoder of the three-hidden family,
$\mathbf{z}=\mathbf{W}_e^{(2)}\sigma(\mathbf{W}_e^{(1)}\mathbf{x}+\mathbf{b}_e^{(1)})+
\mathbf{b}_e^{(2)}$.  The assignment statement calls this the ``2-hidden'' autoencoder, counting
only the layers beyond the bottleneck that carry a nonlinearity; we call it three-hidden throughout
because the object actually built has three hidden layers of widths $400$, $k$ and $400$, while the
one-hidden object of Task 3 has exactly one, of width $k$.  Both names refer to the same two models,
and the naming is stated so that the count never drifts between the code, this report and the
figures.

\begin{table}[htbp]\centering
\caption{Task 4: validation accuracy (\%) of each classifier on the frozen three-hidden code at
each width.  Blank cells are degenerate runs excluded from selection by the $25\,\%$ rule of
Section~\ref{sec:arch:opt}.}\label{tab:task4}
\begin{tabular}{@{}lrrrr@{}}
\toprule Architecture & $k=32$ & $k=64$ & $k=128$ & $k=256$ \\ \midrule
3L\_A & \RESULT & \RESULT & \RESULT & \RESULT \\
3L\_B & \RESULT & \RESULT & \RESULT & \RESULT \\
4L\_A & \RESULT & \RESULT & \RESULT & \RESULT \\
5L\_A & \RESULT & \RESULT & \RESULT & \RESULT \\
\bottomrule
\end{tabular}
\end{table}

%%% RESULT: at $k=256$ the three-hidden code produced chance-level (20.00%) collapses for three of
%%% the four classifiers.  Name the three, state that the survivor was selected and that its test
%%% accuracy was recomputed for it, and report this as an honest negative finding about sigmoid
%%% depth rather than as a defect of the pipeline.

\paragraph{Depth does not buy accuracy.}
Figure~\ref{fig:task4:acc} is the direct counterpart of Figure~\ref{fig:task3:acc}: both families
emit a $\R^{k}$ bottleneck at the same $k$, the classifiers and the protocol are identical, and the
difference is the two $400$-wide sigmoid layers in this encoder.  The expected reading is that
accuracy at matched $k$ is \emph{not} better than Task 3's and is plausibly worse, and that the
widths at which it is worse follow the pattern of the collapses recorded in Table~\ref{tab:task4}.
The mechanistic reading, developed in Section~\ref{sec:obs}, is that a sigmoid stack is a
composition of saturating squashes whose product gradient shrinks with depth, so the encoder's own
training is at risk before the classifier is involved; a code that was never fully trained cannot
classify better than one that was.

\begin{figure}[htbp]\centering
\maybeinc[width=0.84\linewidth]{task4/08_comparisons/task4_accuracy_vs_bottleneck.png}
\caption{Task 4: validation accuracy against bottleneck width $k$ for the frozen three-hidden code.
Read against Figure~\ref{fig:task3:acc} at matched $k$.}\label{fig:task4:acc}
\end{figure}

The best three-hidden configuration is $k=\RESULT$ with the \RESULT{} classifier, at
$\Acc_{\mathrm{val}}=\RESULTPCT$ and $\Acc_{\mathrm{test}}=\RESULTPCT$; its test confusion matrix is
Figure~\ref{fig:task4:cm}, and it shows the same dominant $5\leftrightarrow6$ confusion as
Figures~\ref{fig:task1:cm} and~\ref{fig:task3:cm}.

\begin{figure}[htbp]\centering
\maybeinc[width=0.78\linewidth]{task4/03_confusion_matrices/task4_k128_3L_A_confusion_test.png}
\caption{Task 4: test confusion matrix of the best-by-validation configuration.  Compare with
Figures~\ref{fig:task1:cm} and~\ref{fig:task3:cm}.}\label{fig:task4:cm}
\end{figure}

\section{Task 5: Denoising Autoencoders}\label{sec:task5}
The one-hidden autoencoder is retrained twice with the corruption of Eq.~\eqref{eq:noise} at
$\rho=0.20$ and $\rho=0.40$, targeting the clean image per Eq.~\eqref{eq:dae}.  The bottleneck
width is the one selected in Task 3, $k=\RESULT$, because the assignment fixes it that way and
because it is the only width at which the denoising code is compared against a known plain-code
baseline rather than against itself.  Table~\ref{tab:task5clf} reports the classification
accuracy of both codes, alongside the reconstruction errors.

\begin{table}[htbp]\centering
%%% RESULT: replace every \RESULT with the measured value, and check that the
%%% plain row matches the $k=\RESULT$ row of Table~\ref{tab:task5clf}.
\caption{Task 5: classification accuracy from the frozen denoising bottleneck
representation against the plain one-hidden result at the same $k$ (from
Table~\ref{tab:task3}), together with the post-training reconstruction error of
each model on the test split and the Assignment-3 baseline.}
\label{tab:task5clf}
\begin{tabular}{@{}lrrrr@{}}
\toprule Model & Val acc. (\%) & Test acc. (\%) & Test $E_{\mathrm{rec}}$ & $\Delta$ vs. plain \\
\midrule Plain one-hidden, $k=\RESULT$ & \RESULT & \RESULT & \RESULT & --- \\
Denoising, $\rho=0.20$ & \RESULT & \RESULT & \RESULT & \RESULT \\
Denoising, $\rho=0.40$ & \RESULT & \RESULT & \RESULT & \RESULT \\ \midrule
Assignment 3 baseline, raw input & 98.84 & 98.76 & --- & --- \\
\bottomrule
\end{tabular}
\end{table}

\begin{figure}[htbp]\centering
\maybeinc[width=0.92\linewidth]{task5/04_reconstructions/task5_dae_noise20_triptych_test.png}
\caption{Task 5, $\rho=0.20$: clean original, corrupted input and reconstruction for one image per
class from the test split.  The middle row is the model's actual input; the third row is its
estimate of the top row.  The $\rho=0.40$ triptych is archived in the appendix.}
\label{fig:task5:trip20}
\end{figure}

\paragraph{Reading the triptych.}
Figure~\ref{fig:task5:trip20} answers a question the numbers cannot.  The denoising objective
forces the code to be a sufficient statistic of the \emph{clean} image rather than of the input it
is shown, so a successful model should reconstruct strokes that the mask destroyed: the corrupted
input has holes where the code supplies them.  What to check is whether the reconstruction fills
those holes with structure consistent with the surviving pixels, which is genuine denoising, or
simply emits a class-conditional average, which is what a $\R^{k}$ bottleneck does when $k$ is too
small to encode identity.  The two are distinguishable in the figure: the first preserves the fine
details of the top row wherever the middle row still contains them, while the second converges
every class toward a blurry prototype.  Between the two noise levels the relevant comparison is not
which $\rho$ has the lower error --- the model trained at higher $\rho$ is asked a harder question
--- but how gracefully each degrades.

The expected reading of Table~\ref{tab:task5clf}, developed in Section~\ref{sec:obs}, is that a
modest corruption fraction costs very little classification accuracy while buying a much more robust
reconstruction, whereas a large one eventually pays for its robustness with discriminative accuracy.
That trade-off is the central result of Task 5: the bottleneck that best reconstructs and the
bottleneck that best classifies are not the same object.

\section{Task 6: Weight Visualisation}\label{sec:task6}
Three one-hidden encoders are compared at the width selected in Task 3: the plain autoencoder and
the two denoising ones at $\rho=0.20$ and $\rho=0.40$.  For every unit $j$ of each encoder we
display the maximally activating training image $\mathbf{x}^{*}_j$ and the encoder weight vector
$\mathbf{w}_{e,j}$, both reshaped to $28\times28$, per Eq.~\eqref{eq:maxact}.

\begin{figure}[htbp]\centering
\maybeinc[width=0.92\linewidth]{task6/06_weights_and_activations/01_plain_AE/task6_plain_ae_maxact_and_weight.png}
\caption{Task 6, plain autoencoder: for each unit, the training image that maximally activates it
(left) paired with the incoming encoder weight vector (right), both reshaped to $28\times28$.  Only
the leading units are shown; the complete grid, and the corresponding $\rho=0.20$ and
$\rho=0.40$ grids, are archived in the appendix.}\label{fig:task6:plain}
\end{figure}

\paragraph{What the weight grid means.}
Because the encoder is a single linear map, $\mathbf{w}_{e,j}$ \emph{is} a template: the unit's
pre-activation is the inner product of that template with the image, so the displayed vector is
literally the shape the unit is most sensitive to.  Three readings follow.  A unit whose weight
vector is a coherent stroke-like structure is a stroke detector; a unit whose vector is close to a
blurred class mean is responding to coarse intensity layout rather than to any local feature; and a
unit whose vector is near-zero has effectively opted out of the code, which is the weight-space
signature of a wasted bottleneck unit.  The paired maximally activating image then confirms or denies
the template reading: if $\mathbf{x}^{*}_j$ shows a strong occurrence of the template's structure,
the unit is functioning as advertised, whereas if it shows an outlier the unit is firing on an
atypical combination and carries little general information.

\paragraph{Effect of the noise level.}
Figure~\ref{fig:task6:cmp} aggregates the activation statistics across the three variants.  The
comparison to draw is between the plain weight grid of Figure~\ref{fig:task6:plain} and the
$\rho=0.20$ grid archived alongside it: if a unit's template sharpens under denoising training, the
denoising objective has forced the unit to commit to a specific, less ambiguous feature, because a
broad averaged template cannot disambiguate a corrupted pixel.  Whether that sharpening extends to
$\rho=0.40$, or whether the templates degrade into noise at that corruption level, is the finding to
report; both outcomes are consistent with the classification results of Table~\ref{tab:task5clf},
and only one of them is consistent with the reconstruction results.

\begin{figure}[htbp]\centering
\maybeinc[width=0.86\linewidth]{task6/08_comparisons/task6c_mean_activation_by_variant.png}
\caption{Task 6: mean absolute activation per unit across the plain and the two denoising encoders.
Units whose response collapses under corruption have effectively left the code.}\label{fig:task6:cmp}
\end{figure}

\section{Results Summary and Comparison}\label{sec:results}
Table~\ref{tab:summary} collects the best-by-validation configuration of every task together with
the Assignment-3 baseline.  The comparison is the point of the assignment, and it is a comparison
at unequal batch settings, which is stated rather than hidden: Assignment 3's best model was trained
at batch size $1$, while every model here is trained full-batch, so a small systematic difference
is expected and differences of well under a percentage point should not be over-read.

\begin{table}[htbp]\centering
\caption{Best-by-validation result of every task, compared with the Assignment-3 baseline.
$\Delta_{\mathrm{A3}}$ is the difference from the $98.76\,\%$ test accuracy of the Assignment-3
baseline, in percentage points.}\label{tab:summary}
\begin{tabular}{@{}llrrrr@{}}
\toprule Task & Representation & Width & Val (\%) & Test (\%) & $\Delta_{\mathrm{A3}}$ \\ \midrule
A3 & Raw pixels, A\_128\_64\_32 & 784 & 98.84 & 98.76 & --- \\
1  & PCA code & $k=\RESULT$ & \RESULT & \RESULT & \RESULT \\
3  & One-hidden AE code & $k=\RESULT$ & \RESULT & \RESULT & \RESULT \\
4  & Three-hidden AE code & $k=\RESULT$ & \RESULT & \RESULT & \RESULT \\
5  & Denoising code, $\rho=0.20$ & $k=\RESULT$ & \RESULT & \RESULT & \RESULT \\
5  & Denoising code, $\rho=0.40$ & $k=\RESULT$ & \RESULT & \RESULT & \RESULT \\
\bottomrule
\end{tabular}
\end{table}

\begin{figure}[htbp]\centering
\maybeinc[width=0.90\linewidth]{00_summary/summary_accuracy_vs_dimension.png}
\caption{Cross-task summary: validation and test accuracy against representation width for all
methods on one axis.  This is the single figure that answers the assignment's central question,
namely which representation buys the most accuracy per unit of width.  Per-task bars against the
baseline are archived in the appendix.}\label{fig:summary:acc}
\end{figure}

\paragraph{The three comparisons that matter.}
First, Task 1 against Task 3 at matched $k$ (Figures~\ref{fig:task1:acc} and~\ref{fig:task3:acc}):
a linear unsupervised code against a nonlinear one, both fed to the same four classifiers.  Second,
Task 3 against Task 4 at matched $k$ (Figures~\ref{fig:task3:acc} and~\ref{fig:task4:acc}): whether
the two $400$-wide sigmoid layers are worth their $\approx628\,784$ parameters.  Third, every task
against the raw-input baseline of $98.76\,\%$ (Figure~\ref{fig:summary:acc}): how much accuracy a
compression to $k=256$, a $3\times$ reduction in width, actually costs.  The third is the
assignment's headline question, and the ordering of the curves in Figure~\ref{fig:summary:acc} is
its answer.

\section{Observations and Analysis}\label{sec:obs}
\paragraph{A learned nonlinear code beats a linear code of the same width.}
The one-hidden autoencoder code at $k=256$ closes most of the gap to the raw-input baseline that the
PCA code at its own best width leaves open.  The mechanism is that PCA is restricted to the $k$
directions of largest \emph{global} variance, and Figure~\ref{fig:task1:variance} shows those
directions saturating quickly: the additional components carry a diminishing share of variance, and
that share is increasingly digit-specific texture rather than class-discriminative structure.  The
autoencoder is under no such restriction, because it may choose a $k$-dimensional basis optimal for
\emph{reconstruction} rather than for variance, and reconstruction of a five-class digit set does
force class information into the code --- a bottleneck cannot reproduce a $5$ from a code that has
discarded everything distinguishing a $5$ from a $6$.

\paragraph{The error structure is inherited, not created.}
Across Tasks 1, 3 and 4 the dominant confusion is the $5\leftrightarrow6$ pair, followed by
$7\leftrightarrow4$ and $5\leftrightarrow7$, in the same order regardless of representation.  This
is the most substantive finding in the report: the reduced representations are not introducing new
failure modes, they are compressing the same information the raw-input classifier used and
inheriting its errors.  The confusions are themselves geometrically sensible --- a $5$ and a $6$
differ mainly in the curvature of the lower bowl, while a $4$ and a $7$ share an upper diagonal and
an upright stem.  A code that preserved those distinctions perfectly would show a different pattern,
so the stability of the pattern is direct evidence about what the codes keep.

\paragraph{Deeper encoders hurt rather than help.}
The three-hidden code does not match the one-hidden code at equal $k$ (Figure~\ref{fig:task4:acc}
against Figure~\ref{fig:task3:acc}), and the chance-level collapses at the widest settings are the
extreme form of the same phenomenon.  The explanation is a property of the architecture rather than
of this dataset: a sigmoid stack is a composition of saturating squashes whose product gradient
shrinks with depth, so the encoder's own reconstruction is at risk before the classifier is
involved, and a code that was never fully optimised cannot classify better than one that was.  The
honest conclusion is that the two $400$-wide sigmoid layers are a cost with no measurable benefit
here --- not that the assignment's prescribed architecture is defective, but that a linear bottleneck
preceded by a deep sigmoid encoder is a harder optimisation problem than the assignment's own
objective requires.

\paragraph{Denoising trades discriminative accuracy for robustness.}
Training against a corrupted input with a clean target changes what the bottleneck must encode.  The
$20\,\%$ model gives up very little classification accuracy (Table~\ref{tab:task5clf}) while
reconstructing from an input that has a fifth of its pixels zeroed (Figure~\ref{fig:task5:trip20});
the $40\,\%$ model is a substantially more robust denoiser but eventually pays for that robustness.
This is the trade-off in plain terms: the bottleneck that best reconstructs is not the bottleneck
that best classifies, and the denoising objective moves the code along that trade-off curve in a
known direction.  It is also why Figure~\ref{fig:task6:cmp} should not be read as the noisiest model
being the most robust model --- robustness in reconstruction and discriminative power are different
quantities, and Task 6's per-unit activation statistics are the level at which the two become
separately visible.

\paragraph{Sigmoid-depth collapse is a reportable negative result.}
Several runs finished at exactly $20.00\,\%$ validation accuracy, which is chance for $C=5$, while
the inherited single-epoch stopping rule would have certified at least one of them as converged.
Those runs are excluded from model selection by an explicit degeneracy flag rather than quietly
dropped, and their confusion matrices are archived as the evidence for the exclusion.  Reporting
them is part of the result: the naive tolerance rule is unsound, and on a plateau that oscillates
upward it is unsound in the direction of false confidence.

\section{Conclusion}\label{sec:conclusion}
We trained PCA and two autoencoder families on a five-class MNIST subset, classified the resulting
compressed representations with a fixed set of four FCNNs, extended the one-hidden family to
denoising variants, and visualised the encoder units of all three.  The best result is the
one-hidden code at $k=256$, which recovers most of the Assignment-3 raw-input accuracy from an
eighth of the width; the best linear result is worse than that, and the deeper encoder family is
worse again.  The reconstruction and denoising results show that the bottleneck width that best
reconstructs and the bottleneck width that best classifies are not the same, and that the denoising
objective moves the code along that trade-off in a controlled way.  Two process conclusions are as
important as the accuracy numbers: the inherited single-epoch stopping rule is unsound and was
replaced by a patience window, and architecture and epoch counts are reported from the live code
rather than from an earlier reading of it.  Future work would replace the sigmoid stack with an
activation that does not saturate asymmetrically, and would test whether the ordering found here
survives on the full ten-class MNIST.

\section{References}\label{sec:refs}
\begin{thebibliography}{9}
\bibitem{goodfellow} I.~Goodfellow, Y.~Bengio and A.~Courville, \emph{Deep Learning}, MIT Press,
  2016.  Chapters 5 and 14 (representation learning and autoencoders).
\bibitem{jolliffe} I.~T.~Jolliffe, \emph{Principal Component Analysis}, 2nd~ed., Springer, 2002.
\bibitem{vincent} P.~Vincent, H.~Larochelle, Y.~Bengio and P.-A.~Manzagol, ``Extracting and
  composing robust features with denoising autoencoders,'' in \emph{Proc.\ 25th International
  Conference on Machine Learning}, 2008.
\bibitem{kingma} D.~P.~Kingma and J.~Ba, ``Adam: A method for stochastic optimization,'' in
  \emph{Proc.\ 3rd International Conference on Learning Representations}, 2015.
\bibitem{glorot} X.~Glorot and Y.~Bengio, ``Understanding the difficulty of training deep
  feedforward neural networks,'' in \emph{Proc.\ 14th International Conference on Artificial
  Intelligence and Statistics}, 2010.
\bibitem{lecun} Y.~LeCun, L.~Bottou, Y.~Bengio and P.~Haffner, ``Gradient-based learning applied
  to document recognition,'' \emph{Proceedings of the IEEE}, 86(11):2278--2324, 1998.
\bibitem{a3} Group~11, \emph{Assignment 3: Fully Connected Neural Networks for Image
  Classification}, CS601T Deep Learning, Indian Institute of Technology Delhi, August--November
  2026.
\end{thebibliography}

\appendix
\section{Figure Gallery}\label{sec:gallery}
This appendix exists so that the question ``where is the image for $X$?'' has a single answer.  It
is additive: every figure is also placed inline at the point in the body text where it is discussed,
and nothing here replaces that placement.  The list below is generated automatically from the
\code{\textbackslash caption} and \code{\textbackslash label} of every figure in the document, so it
cannot drift out of step with the body.

\listoffigures

\subsection{How figures are archived}\label{sec:gallery:policy}
Figures are written by the plotting stage of the pipeline under
\code{Group11\_Assignment4\_code/results\_final/plots/}, the directory declared in this document's
\code{\textbackslash graphicspath}.  The tree is organised by task and then by question, with the
same conceptual sub-folder name reused across all tasks so that a question has exactly one place to
look:
\begin{itemize}[leftmargin=1.6em,itemsep=1pt]
\item \code{plots/task1/01\_training\_curves}, \code{02\_accuracy},
  \code{03\_confusion\_matrices}, \code{07\_representations}, \code{08\_comparisons}, with the same
  numbering scheme in tasks 2--6;
\item \code{plots/task2/04\_reconstructions} and \code{05\_reconstruction\_error};
\item \code{plots/task5/04\_reconstructions} (clean / corrupted / reconstruction triptychs) and
  \code{05\_reconstruction\_error};
\item \code{plots/task6/06\_weights\_and\_activations/}, further split into
  \code{01\_plain\_AE}, \code{02\_denoising\_AE\_20pct} and \code{03\_denoising\_AE\_40pct};
\item \code{plots/00\_summary/} for cross-task comparisons, indexed by
  \code{plots/00\_summary/README.txt}, which records the folder-to-question mapping.
\end{itemize}
File names are self-describing and encode task, width, architecture, noise level and split --- for
example \code{task1\_k32\_3L\_A\_confusion\_test.png} --- so that figures from different runs cannot
overwrite one another.

\subsection{To regenerate and re-embed a figure}\label{sec:gallery:regen}
Every \code{\textbackslash maybeinc} call in this document names the exact path it expects, relative
to \code{plots/}.  When that file is absent at compile time a labelled placeholder frame is drawn
instead, which is why this document compiles with an empty results tree just as well as with a
complete one.  To embed a figure verbatim once it has been archived, replace the guarded call
\begin{lstlisting}[language={},basicstyle=\small\ttfamily,numbers=none,
caption={Guarded include used throughout this report.},label={lst:maybeinc}]
\maybeinc[width=0.84\linewidth]{task1/07_representations/task1_variance_retained.png}
\end{lstlisting}
with the plain include
\begin{lstlisting}[language={},basicstyle=\small\ttfamily,numbers=none,
caption={Verbatim include, to be used once the file exists.},label={lst:rawinc}]
\includegraphics[width=0.84\linewidth]{task1/07_representations/task1_variance_retained.png}
\end{lstlisting}
Both forms resolve through the same \code{\textbackslash graphicspath}, so no path changes elsewhere
in the document.

%%% RESULT: after the full pipeline has been run, list the count of archived PNG files per task
%%% folder here, and confirm that the figure labels above match the set of figures referenced in
%%% the body text.  Any figure generated by the pipeline but not referenced should either be
%%% referenced or explicitly noted as an extra diagnostic.

\end{document}
